\section{La función \texorpdfstring{$F$}{F} y el valor óptimo}\label{sec:F}

Los cuatro factores $r_R,r_C,u_D,u_A$ guardan información sobre cada posición. Al promediar
el producto de las dos vacancias diagonales con el peso $r_R(r)\,r_C(c)$ se
obtiene un número, $F(x)$, que depende sólo de la fracción de filas vacías $x$ y que permite
recuperar, integrando su logaritmo, el valor óptimo de la entropía en masa uno.

Para $x\in(0,1)$ y $\rho=1-x$, la función $F$ está definida en~\eqref{eq:F}. Integrando
$\nu=b\,r_Rr_Cu_Du_A$, con $b=\lambda x^2$ e $\int\nu=1-x$, se obtiene
\begin{equation}\label{eq:F_integral}
 \int_{[0,1]^2}\Theta_x\,u_D(r-c)\,u_A(r+c)\,dr\,dc=\frac{1-x}{b}=F(x),\qquad \Theta_x=r_Rr_C .
\end{equation}
El peso $\Theta_x$ integra uno. La igualdad es una consecuencia algebraica de la ecuación de Euler.

\begin{teorema}[Analiticidad de $F$]\label{thm:F}
La función $F$ se extiende a una función analítica real y positiva en un entorno de $[0,1]$, con
$F(0)=1/b_0$ y $F(1)=1$.
\end{teorema}

\begin{proof}
Cerca de $x=0$, $F(x)=(1-x)/b_x$ con $b_x=e^{\gamma(x)}$ analítica (Teorema~\ref{thm:rama}), y
$F(0)=1/b_0>0$. En $0<x<1$, $\lambda$ es analítica en $\rho$ (Proposición~\ref{prop:interior}), y
\eqref{eq:F} da la analiticidad de $F$. Cerca de $x=1$, la Proposición~\ref{prop:diluido} da
$\lambda(1-x)=(1-x)A(1-x)$, de modo que
\[
 F(x)=\frac1{A(1-x)\,x^2},\qquad F(1)=1 .
\]
Las tres expresiones coinciden donde sus dominios se solapan, porque todas proceden de los mismos
maximizadores. Como $F>0$ en $(0,1)$, $F(0)>0$ y $F(1)=1$, por compacidad $F$ es positiva en un
entorno de $[0,1]$, reduciéndolo si hace falta.
\end{proof}

\begin{proposicion}[El valor óptimo en masa uno]\label{prop:intlogF}
\[
 \int_0^1\log F(x)\,dx=1+\sopt(1).
\]
\end{proposicion}

\begin{proof}
Para $x\in(0,1)$ y $\rho=1-x$, el Teorema~\ref{thm:max}(d) da $\lambda=e^{-\sopt{}'(\rho)}$, de modo que
\[
 \log F(1-\rho)=\log\rho+\sopt{}'(\rho)-2\log(1-\rho).
\]
Además $\sopt(\rho)\to0$ cuando $\rho\downarrow0$. En efecto, $G(m)\le m$ en $[0,1]$, porque
$-\log(1-m)\le m/(1-m)$; por Jensen y el Lema~\ref{lem:fubini},
$s[\nu]\le-3\rho-\rho\log\rho+\sum_f\int m_f=\rho-\rho\log\rho$ en $\cA_\rho$, y la densidad constante da
$\sopt(\rho)\ge-3\rho-\rho\log\rho$, porque $G\ge0$. Por la Proposición~\ref{prop:masa_uno},
$\sopt(\rho)\to\sopt(1)$ cuando $\rho\uparrow1$. Como $\log F$ es acotada en $[0,1]$
(Teorema~\ref{thm:F}), se integra en $[\varepsilon,1-\varepsilon]$ y se hace $\varepsilon\to0$: los
términos $\log\rho$ y $-2\log(1-\rho)$ aportan $-1$ y $2$, y
$\int_\varepsilon^{1-\varepsilon}\sopt{}'=\sopt(1-\varepsilon)-\sopt(\varepsilon)\to\sopt(1)$.
\end{proof}

\begin{proposicion}[Expansión en el tablero lleno]\label{prop:borde}
Cuando $x\downarrow0$,
\[
 \sopt(1-x)-\sopt(1)=2x\log\frac1x+O(x).
\]
\end{proposicion}

\begin{proof}
Para $\rho=1-y$ con $0<y\le\frac12$, $\sopt{}'(\rho)=-\log\lambda=2\log y-\log b_y$, y
$|\log b_y|\le\max\{\log2,\log b_+\}$ por la Proposición~\ref{prop:cotas}. Como
$\sopt(\rho)\to\sopt(1)$ (Proposición~\ref{prop:masa_uno}) y $\log y$ es integrable,
\[
 \sopt(1)-\sopt(1-x)=\int_0^x\bigl(2\log y-\log b_y\bigr)\,dy=2x\log x-2x+O(x),
\]
que es la expansión enunciada.
\end{proof}

\begin{observacion}\label{obs:borde_fino}
Como $b$ es analítica en $x=0$ (Teorema~\ref{thm:rama}), $\log b_y=\log b_0+O(y)$, y la misma cuenta da
$\sopt(1-x)-\sopt(1)=2x\log\frac1x+\bigl(2-\log F(0)\bigr)x+O(x^2)$, porque $F(0)=1/b_0$.
\end{observacion}

Con el Teorema~\ref{thm:F} y las Proposiciones~\ref{prop:intlogF} y~\ref{prop:borde} queda probado el
Teorema~\ref{thm:F_intro}.
